\documentclass[11pt]{article}
\usepackage{amsmath, amssymb, geometry, hyperref, graphicx}
\geometry{a4paper, margin=1in}

\title{Work Stream 2: Deriving the Thermal Energy Equation and Vaporization Limit}
\author{MechanicaFluidorum Program \& Community Contributors (esp. \texttt{u/vibe0009} \& \texttt{u/Little-Name9809})}
\date{September 2026}

\begin{document}
\maketitle

\section{Introduction}
A central question raised by \texttt{u/vibe0009} in \texttt{r/FluidMechanics} was how enstrophy divergence implies thermal shock when the incompressible Navier-Stokes equations lack a temperature equation. To address this, we must step outside the idealized incompressible frame and couple the viscous dissipation to the thermal energy equation.

\section{Thermal Coupling}
The complete energy equation for a fluid includes the viscous dissipation function $\Phi$:
\begin{equation}
    \rho c_p \left( \frac{\partial T}{\partial t} + u \cdot \nabla T \right) = k \nabla^2 T + 2\mu \, \mathbf{D}(u) : \mathbf{D}(u)
\end{equation}
The contraction of the strain rate tensor $\mathbf{D}(u)$ is directly proportional to the enstrophy density $|\nabla \times u|^2$ in incompressible flow.

\section{Heating by Pure Force}
\texttt{u/Little-Name9809} asked whether this suggests a mechanism for vaporizing a fluid using only mechanical force (no external heat). In the OpenAI construction the \emph{local} enstrophy density scales as $|\omega|^2\sim\tau^{-2-2h}=\tau^{-2.010}$ (the \emph{integrated} enstrophy $\Omega=\int|\omega|^2dx$, which also absorbs the shrinking volume, scales as $\tau^{-0.515}$; $h=1/200$), so the local heat generation rate $\varepsilon=\nu|\omega|^2$ diverges.

The temperature rise follows in three lines. The core keeps a radial Reynolds number of order one, so its radius is the diffusive length $\ell_r^2=\nu t$, where $t$ is the physical time remaining. The dissipation per unit mass is $\varepsilon=\nu(u/\ell_r)^2$. Acting over the time $t$,
\begin{equation}
    \Delta T \;\simeq\; \frac{\varepsilon\,t}{c_p} \;=\; \frac{\nu u^2 t}{\ell_r^2\,c_p} \;=\; \frac{u^2}{c_p}\;\sim\;\tau^{-1.010},
\end{equation}
an Eckert number of order one: the kinetic energy per unit mass of the core is converted to heat within the core on its own collapse time. For water ($c_p=4184$\,J/kg/K) this gives $\Delta T\approx48$\,K at $\mathrm{Ma}=0.3$ and $\approx540$\,K at $\mathrm{Ma}=1$; boiling ($\Delta T=73$\,K) is reached at $u\approx553$\,m/s, $\mathrm{Ma}\approx0.37$, about $3$--$4.5$\,ps before the mathematical blow-up (an earlier version of this note gave $45.6$\,ps, from an illustrative constant that has since been removed). Heating is therefore an order-one effect at the same scale, $\ell_*=\nu/c$, at which compressibility and rarefaction appear; it is not an independent, earlier limit. A plasma ($\sim10^4$\,K) would need $\mathrm{Ma}\approx4$, far outside the incompressible model, so ``plasma by pure force'' is not supported by this estimate.

\emph{Caveat.} The estimate assumes the heat stays in the core over the collapse time. The thermal diffusion length $\sqrt{\alpha t}$ compares with $\ell_r=\sqrt{\nu t}$ as $\mathrm{Pr}^{-1/2}$: for water ($\mathrm{Pr}\approx7$) conduction removes only a fraction of the heat, so the estimate stands; for gases ($\mathrm{Pr}\approx0.7$) it is an upper bound. Either way the conclusion is that the constant-temperature, constant-density assumptions of the model stop holding at $\ell_*$, not that any particular phase change is predicted.

\section{Acknowledgments}
Thank you to \texttt{u/vibe0009} and \texttt{u/Little-Name9809} for prompting this explicit thermal derivation. \textit{Rendons \`a C\'esar ce qui appartient \`a C\'esar.}

\end{document}
